% ============================================================================
% Fichier  : partie5/P05_C15_mesures_organisationnelles.tex
% Titre    : Mesures Organisationnelles et Humaines
% Auteur   : MINKA MI NGUIDJOÏ Thierry Emmanuel
% ============================================================================
% Sources : ISO/IEC 27002:2022 (thèmes organisationnel et personnes). Fictif.
% ============================================================================

\chapter{Mesures Organisationnelles et Humaines}
\label{chap:orga}
\chaptermark{Mesures organisationnelles}

\epigraph{%
    « La technique protège les machines~;\\
    la sensibilisation protège les décisions des gens. »%
}{Principe du cours}

\niveauChapitre{Intermédiaire}{La partie IV.}

\begin{boxobjectifs}
\begin{itemize}
    \item Identifier les mesures organisationnelles clés.
    \item Comprendre le poids du facteur humain dans la sécurité.
    \item Concevoir une campagne de sensibilisation.
    \item Sécuriser le cycle de vie de l'agent (arrivée, mobilité, départ).
    \item Encadrer les fournisseurs et le télétravail.
\end{itemize}
\end{boxobjectifs}

\bigskip

% ============================================================================
\section*{Question directrice : comment agir sur ce qui n'est ni un logiciel ni un câble ?}
% ============================================================================

L'incident de la mairie est né d'un geste humain~: brancher une clé. Les
mesures techniques ne suffisent pas~: il faut agir sur l'organisation et sur les
personnes. Ce chapitre traite ces mesures.

\medskip
\textbf{Plan du chapitre.}
\begin{enumerate}
    \item Les mesures organisationnelles (section~\ref{sec:c15-orga}).
    \item Le facteur humain (section~\ref{sec:c15-humain}).
    \item La sensibilisation (section~\ref{sec:c15-sensibilisation}).
    \item Le cycle de vie de l'agent (section~\ref{sec:c15-cyclevie}).
    \item Fournisseurs et télétravail (section~\ref{sec:c15-tiers}).
\end{enumerate}

\separateur

% ============================================================================
\section{Les Mesures Organisationnelles}
\label{sec:c15-orga}
% ============================================================================

\begin{tableaucours}{L{4.2cm}L{8.0cm}}{%
    \tableauHeader{Mesure} &
    \tableauHeader{Ce qu'elle installe}
}
Politiques et procédures & Le cadre écrit des règles \\
Rôles et responsabilités & Qui fait quoi (matrice RACI) \\
Classification de l'information & Des niveaux et des règles de manipulation \\
Gestion des incidents & Un processus de signalement et de traitement \\
Relations avec les fournisseurs & Des exigences de sécurité contractuelles \\
\end{tableaucours}
\vspace{-0.8em}
\begin{center}{\small\itshape Tableau~15.1. Mesures organisationnelles (ISO/IEC 27002)}\end{center}

Ces mesures ne coûtent souvent que du temps et de la rigueur~: elles sont les
premières à déployer, avant tout investissement technique.

% ============================================================================
\section{Le Facteur Humain}
\label{sec:c15-humain}
% ============================================================================

\begin{boxdef}
Le \termecle{facteur humain}~: part des incidents liée aux comportements des
personnes (erreur, négligence, manipulation). Il est, dans la plupart des
études, la première cause d'incident.
\end{boxdef}

L'agent n'est pas le maillon faible par nature~: il le devient faute de
formation et de règles claires. Sensibiliser transforme le maillon faible en
première ligne de défense.

\begin{boxscenario}{Opération MAKAK. La clé piégée}
\textbf{Constatation.} L'agent qui a branché la clé ignorait le risque~: il
n'avait reçu ni règle écrite, ni formation.

\medskip
\textbf{Analyse} (confiance élevée). Le reproche ne vise pas l'agent mais
l'organisation~: sans règle ni sensibilisation, le geste était prévisible.
La mesure corrective est double~: interdire les clés personnelles \emph{et}
former les agents.
\end{boxscenario}

% ============================================================================
\section{La Sensibilisation}
\label{sec:c15-sensibilisation}
% ============================================================================

\begin{tableaucours}{L{4.0cm}L{8.2cm}}{%
    \tableauHeader{Thème} &
    \tableauHeader{Message clé}
}
Mots de passe & Individuels, robustes, jamais partagés \\
Hameçonnage & Vérifier l'expéditeur, ne pas cliquer dans le doute \\
Supports amovibles & N'utiliser que des supports maîtrisés \\
Données personnelles & Ne traiter que le nécessaire, protéger \\
Signalement & Prévenir vite en cas de doute, sans crainte \\
\end{tableaucours}
\vspace{-0.8em}
\begin{center}{\small\itshape Tableau~15.2. Thèmes d'une campagne de sensibilisation}\end{center}

Une campagne efficace est régulière, concrète et bienveillante~: on apprend à
signaler sans être puni. Hameçonnage simulé, affiches et courtes sessions
valent mieux qu'un règlement illisible.

\begin{boxmetier}{Responsable conformité}
\textbf{Prouver la sensibilisation.} La formation doit être tracée~: feuilles
d'émargement, attestations, résultats des tests. En audit, « nous sensibilisons
nos agents » ne vaut rien sans preuve.
\end{boxmetier}

% ============================================================================
\section{Le Cycle de Vie de l'Agent}
\label{sec:c15-cyclevie}
% ============================================================================

La sécurité suit l'agent de son arrivée à son départ~: à l'\termecle{arrivée},
charte signée et accès minimal~; en cas de \termecle{mobilité}, droits
réajustés~; au \termecle{départ}, accès révoqués et matériel restitué. Un
compte d'agent parti et toujours actif est une porte ouverte.

% ============================================================================
\section{Fournisseurs et Télétravail}
\label{sec:c15-tiers}
% ============================================================================

\begin{boxprocedure}
\textbf{Constatation.} L'éditeur du logiciel d'état civil accède à distance aux
données de la mairie, sans clause de sécurité au contrat.

\medskip
\textbf{Analyse.} C'est un risque de tiers~: la mairie reste responsable des
données que son prestataire manipule. La mesure impose des \emph{exigences de
sécurité contractuelles} et un contrôle des accès distants. Le télétravail des
agents suit la même logique~: règles écrites et moyens maîtrisés.
\end{boxprocedure}

\separateur

\begin{boxresume}
\textbf{Ce qu'il faut retenir de ce chapitre~:}
\begin{itemize}
    \item Les \textbf{mesures organisationnelles} posent le cadre écrit, à
          faible coût.
    \item Le \textbf{facteur humain} est la première cause d'incident~; la
          formation en fait une défense.
    \item Une \textbf{campagne de sensibilisation} est régulière, concrète,
          tracée.
    \item La sécurité suit le \textbf{cycle de vie de l'agent}, du recrutement
          au départ.
    \item \textbf{Fournisseurs} et \textbf{télétravail} exigent des règles et
          des contrats.
\end{itemize}
\end{boxresume}

\begin{boxinfo}
\textbf{Points de contrôle de l'élève.} Je sais~: citer des mesures
organisationnelles~; expliquer le facteur humain~; concevoir une campagne~;
sécuriser le cycle de vie de l'agent~; encadrer fournisseurs et télétravail.
\end{boxinfo}

\bigskip

\begin{boxexo}[title={\ding{45}\quad Exercice~15.1. Mesures organisationnelles \niveaubase}]
Citez trois mesures organisationnelles de l'ISO/IEC 27002.
\end{boxexo}

\begin{boxexo}[title={\ding{45}\quad Exercice~15.2. Facteur humain \niveaubase}]
Pourquoi dit-on que l'humain est la première cause d'incident~?
\end{boxexo}

\begin{boxexo}[title={\ding{45}\quad Exercice~15.3. Thèmes \niveaubase}]
Citez quatre thèmes d'une campagne de sensibilisation.
\end{boxexo}

\begin{boxexo}[title={\ding{45}\quad Exercice~15.4. Blâmer qui \niveaumoyen}]
Faut-il blâmer l'agent qui a branché la clé~? Argumentez.
\end{boxexo}

\begin{boxexo}[title={\ding{45}\quad Exercice~15.5. Prouver la formation \niveaumoyen}]
Comment la mairie prouve-t-elle, en audit, qu'elle sensibilise ses agents~?
\end{boxexo}

\begin{boxexo}[title={\ding{45}\quad Exercice~15.6. Départ d'un agent \niveaumoyen}]
Quelles actions de sécurité au départ d'un agent~?
\end{boxexo}

\begin{boxexo}[title={\ding{45}\quad Exercice~15.7. Risque de tiers \niveaumoyen}]
Pourquoi la mairie reste-t-elle responsable des données manipulées par son
prestataire~?
\end{boxexo}

\begin{boxexo}[title={\ding{45}\quad Exercice~15.8. Télétravail \niveaumoyen}]
Citez trois règles de sécurité pour le télétravail d'un agent de la mairie.
\end{boxexo}

\begin{boxexo}[title={\ding{45}\quad Exercice~15.9. Plan de campagne \niveauexpert}]
Concevez un plan de campagne de sensibilisation sur un an pour la mairie
d'Eséka.
\end{boxexo}

\begin{boxexo}[title={\ding{45}\quad Exercice~15.10. Clause fournisseur \niveauexpert}]
Rédigez une clause de sécurité à insérer au contrat de l'éditeur du logiciel
d'état civil.
\end{boxexo}

\bigskip

\begin{boxinfo}
\textbf{Pour aller plus loin.}
\begin{itemize}
    \item ISO/IEC 27002, thèmes organisationnel et personnes~\cite{iso27002}~;
          Guide d'hygiène ANSSI~\cite{anssihygiene2017}.
\end{itemize}
\end{boxinfo}

\begin{boxinfo}
\textbf{Transition.} Les mesures humaines posées, le Chapitre~\ref{chap:technique}
traite les mesures techniques~: accès, réseau, chiffrement, journalisation.
\end{boxinfo}
